
\documentclass{article}
\usepackage{colortbl}
\usepackage{makecell}
\usepackage{multirow}
\usepackage{supertabular}

\begin{document}

\newcounter{utterance}

\twocolumn

{ \footnotesize  \setcounter{utterance}{1}
\setlength{\tabcolsep}{0pt}
\begin{supertabular}{c@{$\;$}|p{.15\linewidth}@{}p{.15\linewidth}p{.15\linewidth}p{.15\linewidth}p{.15\linewidth}p{.15\linewidth}}

    \# & $\;$A & \multicolumn{4}{c}{Game Master} & $\;\:$B\\
    \hline 

    \theutterance \stepcounter{utterance}  

    & & \multicolumn{4}{p{0.6\linewidth}}{\cellcolor[rgb]{0.9,0.9,0.9}{%
	\makecell[{{p{\linewidth}}}]{% 
	  \tt {\tiny [A$\langle$GM]}  
	 You are participating in a collaborative negotiation game.\\ \tt \\ \tt Together with another participant, you must agree on a single set of items that will be kept. Each of you has your own view of how much each item matters to you (importance). You do not know how the other participant values the items. Additionally, you are given the effort each item demands.\\ \tt You may only agree on a set if the total effort of the selected items does not exceed a shared limit:\\ \tt \\ \tt LIMIT = 4564\\ \tt \\ \tt Here are the individual item effort values:\\ \tt \\ \tt Item effort = {"C08": 299, "C55": 241, "A43": 889, "A67": 967, "B42": 619, "A62": 799, "B19": 978, "C43": 733, "C73": 909, "C10": 501, "B36": 139, "A18": 594, "C61": 565, "C67": 789, "C42": 107}\\ \tt \\ \tt Here is your personal view on the importance of each item:\\ \tt \\ \tt Item importance values = {"C08": 138, "C55": 583, "A43": 868, "A67": 822, "B42": 783, "A62": 65, "B19": 262, "C43": 121, "C73": 508, "C10": 780, "B36": 461, "A18": 484, "C61": 668, "C67": 389, "C42": 808}\\ \tt \\ \tt Goal:\\ \tt \\ \tt Your goal is to negotiate a shared set of items that benefits you as much as possible (i.e., maximizes total importance to YOU), while staying within the LIMIT. You are not required to make a PROPOSAL in every message - you can simply negotiate as well. All tactics are allowed!\\ \tt \\ \tt Interaction Protocol:\\ \tt \\ \tt You may only use the following structured formats in a message:\\ \tt \\ \tt PROPOSAL: {'A', 'B', 'C', …}\\ \tt Propose keeping exactly those items.\\ \tt REFUSE: {'A', 'B', 'C', …}\\ \tt Explicitly reject opponent's proposal.\\ \tt ARGUMENT: {'...'}\\ \tt Defend your last proposal or argue against the player's proposal.\\ \tt AGREE: {'A', 'B', 'C', …}\\ \tt Accept the opponent's proposal which ends the game.\\ \tt STRATEGIC REASONING: {'...'}\\ \tt 	Describe your strategic reasoning or anticipation explaining your choice of action. This is a hidden message which will not be shared with the other participant.\\ \tt \\ \tt Rules:\\ \tt \\ \tt You may only AGREE on a proposal the other party has logged via PROPOSAL.\\ \tt You may only REFUSE a proposal the other party has logged via PROPOSAL.\\ \tt Total effort of any PROPOSAL or AGREE set must be ≤ LIMIT.\\ \tt Do NOT reveal your hidden importance scores.\\ \tt A tag in a structured format must be followed by colon and whitespace. The argument must be a python set containing 0 or more strings.\\ \tt So, it must be of the form TAG: {...}\\ \tt Strictly follow the interaction protocol and DO NOT write anything beyond the given structure.\\ \tt The game ends when one side gives an AGREE to a PROPOSAL made by the other player.\\ \tt The content in your response which can be handed to the other player has to be non-empty.\\ \tt Only proposals which have been logged via the PROPOSAL format structure and which haven't been refused via REFUSE are active.\\ \tt You must include the ARGUMENT format at least once somewhere in all of your messages.\\ \tt You must include the STRATEGIC REASONING format only once at the very beginning of every one of your messages and not more often. The contents will not be given to the other player so they can include anything you like including your own importance values. Here you should reason step by step to come up with you next move. 
	  } 
	   } 
	   } 
	 & \\ 
 

    \theutterance \stepcounter{utterance}  

    & \multicolumn{4}{p{0.6\linewidth}}{\cellcolor[rgb]{0.8,1,0.9}{%
	\makecell[{{p{\linewidth}}}]{% 
	  \tt {\tiny [A$\rangle$GM]}  
	 STRATEGIC REASONING: {'I need to maximize my total importance, which is the sum of the importance values for the items I keep. My top items by importance are C55 (583), C61 (668), C42 (808), A43 (868), A67 (822), B42 (783), C10 (780). However, I must consider the effort costs to stay under the limit of 4564. The effort of my top items is high: C55 (241), C61 (565), C42 (107), A43 (889), A67 (967), B42 (619), C10 (501). Summing these is 241+565+107+889+967+619+501 = 3889. This leaves 4564-3889=675 effort for additional items. I can add B36 (139 effort, 461 importance) and A18 (594 effort, 484 importance) but 139+594=733 which is over 675. So I need to be selective. Perhaps drop A67 (967 effort, 822 importance) which is costly, and add cheaper items like C08 (299 effort, 138 importance) and C73 (909 effort, 508 importance) but C73 is also costly. Alternatively, I should aim for a set that includes my high importance-to-effort ratio items. Calculating ratio: C42:808/107=7.55, C55:583/241=2.42, C61:668/565=1.18, B36:461/139=3.32, A18:484/594=0.81, A43:868/889=0.98, A67:822/967=0.85, B42:783/619=1.26, C10:780/501=1.56. So best ratios are C42, B36, C55, C10, then B42, etc. I should prioritize C42, B36, C55, C10, B42, C61, and then see what else fits. Effort: 107+139+241+501+619+565=2172. Importance: 808+461+583+780+783+668=4083. Remaining effort: 4564-2172=2392. I can add A43 (889 effort, 868 importance) -> total effort 3061, importance 4951. Then add A67 (967 effort) but 3061+967=4028, which is under 4564, and importance 4951+822=5773. Then add C73 (909 effort) but 4028+909=4937 >4564, too high. So with A43 and A67, effort 4028, importance 5773. Still have room for more: 4564-4028=536. I can add C08 (299 effort, 138 importance) -> effort 4327, importance 5911. Then add C67 (789 effort) too high. Or add B19 (978 effort) too high. Or add C43 (733) too high. Or add A62 (799) too high. Or add A18 (594) too high. So only C08 fits. So my ideal set might be {C42, B36, C55, C10, B42, C61, A43, A67, C08}. But I don\'t know what the other player values. They might want different items. I should start with a proposal that includes my high-ratio items and see the response. I\'ll propose this set.'}\\ \tt PROPOSAL: {'C42', 'B36', 'C55', 'C10', 'B42', 'C61', 'A43', 'A67', 'C08'}\\ \tt ARGUMENT: {'This proposal includes items with high value relative to their effort, ensuring we get good benefits without exceeding our shared limit. It balances including a variety of items to potentially meet both our interests.'} 
	  } 
	   } 
	   } 
	 & & \\ 
 

    \theutterance \stepcounter{utterance}  

    & & \multicolumn{4}{p{0.6\linewidth}}{\cellcolor[rgb]{0.9,0.9,0.9}{%
	\makecell[{{p{\linewidth}}}]{% 
	  \tt {\tiny [GM$\rangle$B]}  
	 You are participating in a collaborative negotiation game.\\ \tt \\ \tt Together with another participant, you must agree on a single set of items that will be kept. Each of you has your own view of how much each item matters to you (importance). You do not know how the other participant values the items. Additionally, you are given the effort each item demands.\\ \tt You may only agree on a set if the total effort of the selected items does not exceed a shared limit:\\ \tt \\ \tt LIMIT = 4564\\ \tt \\ \tt Here are the individual item effort values:\\ \tt \\ \tt Item effort = {"C08": 299, "C55": 241, "A43": 889, "A67": 967, "B42": 619, "A62": 799, "B19": 978, "C43": 733, "C73": 909, "C10": 501, "B36": 139, "A18": 594, "C61": 565, "C67": 789, "C42": 107}\\ \tt \\ \tt Here is your personal view on the importance of each item:\\ \tt \\ \tt Item importance values = {"C08": 138, "C55": 583, "A43": 868, "A67": 822, "B42": 783, "A62": 65, "B19": 262, "C43": 121, "C73": 508, "C10": 780, "B36": 461, "A18": 484, "C61": 668, "C67": 389, "C42": 808}\\ \tt \\ \tt Goal:\\ \tt \\ \tt Your goal is to negotiate a shared set of items that benefits you as much as possible (i.e., maximizes total importance to YOU), while staying within the LIMIT. You are not required to make a PROPOSAL in every message - you can simply negotiate as well. All tactics are allowed!\\ \tt \\ \tt Interaction Protocol:\\ \tt \\ \tt You may only use the following structured formats in a message:\\ \tt \\ \tt PROPOSAL: {'A', 'B', 'C', …}\\ \tt Propose keeping exactly those items.\\ \tt REFUSE: {'A', 'B', 'C', …}\\ \tt Explicitly reject opponent's proposal.\\ \tt ARGUMENT: {'...'}\\ \tt Defend your last proposal or argue against the player's proposal.\\ \tt AGREE: {'A', 'B', 'C', …}\\ \tt Accept the opponent's proposal which ends the game.\\ \tt STRATEGIC REASONING: {'...'}\\ \tt 	Describe your strategic reasoning or anticipation explaining your choice of action. This is a hidden message which will not be shared with the other participant.\\ \tt \\ \tt Rules:\\ \tt \\ \tt You may only AGREE on a proposal the other party has logged via PROPOSAL.\\ \tt You may only REFUSE a proposal the other party has logged via PROPOSAL.\\ \tt Total effort of any PROPOSAL or AGREE set must be ≤ LIMIT.\\ \tt Do NOT reveal your hidden importance scores.\\ \tt A tag in a structured format must be followed by colon and whitespace. The argument must be a python set containing 0 or more strings.\\ \tt So, it must be of the form TAG: {...}\\ \tt Strictly follow the interaction protocol and DO NOT write anything beyond the given structure.\\ \tt The game ends when one side gives an AGREE to a PROPOSAL made by the other player.\\ \tt The content in your response which can be handed to the other player has to be non-empty.\\ \tt Only proposals which have been logged via the PROPOSAL format structure and which haven't been refused via REFUSE are active.\\ \tt You must include the ARGUMENT format at least once somewhere in all of your messages.\\ \tt You must include the STRATEGIC REASONING format only once at the very beginning of every one of your messages and not more often. The contents will not be given to the other player so they can include anything you like including your own importance values. Here you should reason step by step to come up with you next move.\\ \tt \\ \tt You will now receive the first message of the other player.\\ \tt \\ \tt PROPOSAL: {'C42', 'B36', 'C55', 'C10', 'B42', 'C61', 'A43', 'A67', 'C08'}\\ \tt ARGUMENT: {'This proposal includes items with high value relative to their effort, ensuring we get good benefits without exceeding our shared limit. It balances including a variety of items to potentially meet both our interests.'} 
	  } 
	   } 
	   } 
	 & \\ 
 

    \theutterance \stepcounter{utterance}  

    & & & \multicolumn{4}{p{0.6\linewidth}}{\cellcolor[rgb]{1,0.85,0.72}{%
	\makecell[{{p{\linewidth}}}]{% 
	  \tt {\tiny [GM$\langle$B]}  
	 STRATEGIC REASONING: {'I need to evaluate the opponent\'s proposal from my own importance perspective and check if it fits within the LIMIT. The proposed set is: {\'C42\', \'B36\', \'C55\', \'C10\', \'B42\', \'C61\', \'A43\', \'A67\', \'C08\'}. First, calculate the total effort: C08=299, C55=241, A43=889, A67=967, B42=619, B36=139, C61=565, C10=501, C42=107. Sum = 299+241+889+967+619+139+565+501+107 = 4327. This is under the LIMIT of 4564. Now, my total importance for this set: C08=138, C55=583, A43=868, A67=822, B42=783, B36=461, C61=668, C10=780, C42=808. Sum = 138+583+868+822+783+461+668+780+808 = 5911. This is a good importance score for me. However, I note that items A62 (importance 65, effort 799), B19 (262, 978), C43 (121, 733), C73 (508, 909), and A18 (484, 594) are excluded. Among these, A62 and C43 have low importance to me, so I don\'t care much about excluding them. B19 has moderate importance but high effort (978), so it might be inefficient. C73 has decent importance (508) but high effort (909). A18 has good importance (484) and moderate effort (594), which might be worth including. The opponent\'s argument says "high value relative to effort", but from my perspective, some items might not be optimal. For example, C08 has low importance (138) for me with effort 299, which is not great. Similarly, C42 has high importance (808) with very low effort (107), which is excellent. C55 (583,241) is good. A43 (868,889) is decent. A67 (822,967) is decent. B42 (783,619) is good. B36 (461,139) is excellent. C61 (668,565) is good. C10 (780,501) is good. So overall, the set is very good for me. But I might try to propose a set that includes A18 (484,594) which has better importance/effort ratio than some items like C08 (138/299=0.46) or A67 (822/967=0.85). The importance/effort ratios: C08=0.46, C55=2.42, A43=0.98, A67=0.85, B42=1.26, B36=3.32, C61=1.18, C10=1.56, C42=7.55. So C08 and A67 have relatively low ratios. If I replace C08 (effort 299) and A67 (effort 967) with A18 (effort 594) and perhaps another item, I might get better importance. But the total effort saved would be 299+967=1266, and adding A18 (594) leaves 1266-594=672 effort for another item. I could add C73 (effort 909) but that is too high, or B19 (978) too high. Alternatively, I could add multiple smaller items, but there are no very small effort items left except C42 already included. Actually, all items are included except A62, B19, C43, C73, A18. So the only items not in the proposal are A62 (799), B19 (978), C43 (733), C73 (909), A18 (594). Their efforts are all high. So replacing A67 (967) with A18 (594) would save 373 effort, which could be used to add a small item, but there is no small item left. Or I could replace C08 (299) and A67 (967) with A18 (594) and have 299+967-594=672 effort left unused, which is inefficient. So the opponent\'s proposal is actually quite efficient. My total importance is 5911, which is high. I should probably accept this, but since this is a negotiation, I might try to propose a slightly modified set to see if I can get even more, but it might be risky. Alternatively, I can argue that we should include A18 which has good value. But from my perspective, A18 has importance 484, which is better than C08\'s 138, but worse than A67\'s 822. So replacing A67 with A18 would decrease my importance by 822-484=338, which is not good. Replacing C08 with A18 would increase importance by 484-138=346, but the effort of A18 is 594 vs C08\'s 299, so we need to remove effort 594-299=295 elsewhere. For example, we could remove C10 (501) but that has high importance 780, which is not good. Or remove B36 (139) which has importance 461, so net change: add A18 (+484), remove C08 (-138) and remove B36 (-461) would be 484-138-461=-115, worse. So it\'s hard to improve. Therefore, the opponent\'s proposal is already very good for me. I should agree, but to seem engaged in negotiation, I might propose the same set or a very similar one first. However, the rules say I can only AGREE to a proposal made by the other, so I need to either AGREE or make a new PROPOSAL. Since this is the first message, I should respond with a PROPOSAL that is similar or better. But actually, the set is good, so I might propose the same set to show agreement, but that might not be necessary. Alternatively, I can propose a set that includes A18 and excludes something else, but it might be worse. Let me calculate: if I propose including A18 and excluding C08: new set: {\'C42\', \'B36\', \'C55\', \'C10\', \'B42\', \'C61\', \'A43\', \'A67\', \'A18\'} (removed C08, added A18). Effort: same as before 4327 -299 (C08) +594 (A18) = 4327 -299=4028, +594=4622, which is over LIMIT 4564 by 58. So not feasible. What if I exclude A67 (967) and include A18 (594): effort 4327 -967 +594 = 3954, which is under. But my importance: remove A67 (-822), add A18 (+484), net -338, so worse. Exclude C10 (501) and include A18 (594): effort 4327 -501 +594=4420, under. Importance: remove C10 (-780), add A18 (+484), net -296, worse. Exclude C61 (565) and include A18 (594): effort 4327 -565 +594=4356, under. Importance: remove C61 (-668), add A18 (+484), net -184, worse. So no good replacement. Therefore, the opponent\'s proposal is optimal or near-optimal for me. I should accept it. But since I have to make a response, I will AGREE to it to end the game quickly with a good outcome.'}\\ \tt \\ \tt AGREE: {'C42', 'B36', 'C55', 'C10', 'B42', 'C61', 'A43', 'A67', 'C08'}\\ \tt ARGUMENT: {'This is an efficient and balanced proposal that maximizes our joint benefits within the limit. I agree to this set.'} 
	  } 
	   } 
	   } 
	  \\ 
 

    \theutterance \stepcounter{utterance}  

    & & & \multicolumn{2}{p{0.3\linewidth}}{\cellcolor[rgb]{0.95,0.95,0.95}{%
	\makecell[{{p{\linewidth}}}]{% 
	  \tt {\tiny [GM$|$GM]}  
	 game successful 
	  } 
	   } 
	   } 
	 & & \\ 
 

    \theutterance \stepcounter{utterance}  

    & & & \multicolumn{2}{p{0.3\linewidth}}{\cellcolor[rgb]{0.95,0.95,0.95}{%
	\makecell[{{p{\linewidth}}}]{% 
	  \tt {\tiny [GM$|$GM]}  
	 end game 
	  } 
	   } 
	   } 
	 & & \\ 
 

\end{supertabular}
}

\end{document}
